\documentclass[11pt,a4paper]{article}
\usepackage[utf8]{inputenc}
\usepackage[T1]{fontenc}
\usepackage[margin=2.3cm]{geometry}
\usepackage{amsmath,amssymb,amsthm}
\usepackage{booktabs}
\usepackage{hyperref}
\usepackage{enumitem}
\newtheorem{theorem}{Théorème}
\newtheorem{proposition}[theorem]{Proposition}
\newtheorem{lemma}[theorem]{Lemme}
\theoremstyle{definition}
\newtheorem{definition}[theorem]{Définition}
\theoremstyle{remark}
\newtheorem{remark}[theorem]{Remarque}
\newcommand{\F}{\mathbb{F}}
\newcommand{\Z}{\mathbb{Z}}
\newcommand{\Fix}{\operatorname{Fix}}
\newcommand{\rg}{\operatorname{rang}}
\newcommand{\srg}{\operatorname{srg}}
\newcommand{\Aut}{\operatorname{Aut}}
\title{Le graphe de Conway et la branche « involution »\\ \large Notes de recherche : ce que j'explore, ce que j'ai trouvé, ce qui tourne}
\author{Journal de la session du 21 septembre 2026 (Claude, sous la direction de D.~Djematene)}
\date{Version de travail --- tout énoncé est soit cité à une source primaire, soit démontré ici. Révisée le 22 septembre 2026 (voir \texttt{notes/AUDIT\_2026-09-22.md})}
\begin{document}
\maketitle

\section{Le problème, en partant de zéro}
Un \emph{graphe} est un ensemble de sommets reliés par des arêtes. Il est \emph{fortement régulier} de paramètres $(v,k,\lambda,\mu)$ s'il a $v$ sommets,
si chaque sommet a exactement $k$ voisins, si deux sommets voisins ont exactement $\lambda$ voisins communs et deux sommets non voisins exactement $\mu$.
Exemples : le pentagone est un $\srg(5,2,0,1)$ ; le graphe de la tour $3\times3$ (cases d'un échiquier $3\times 3$, reliées si même ligne ou même colonne)
est un $\srg(9,4,1,2)$ --- c'est notre \emph{graphe de contrôle} tout au long de ce travail.

Le \emph{graphe de Conway} serait un $\srg(99,14,1,2)$ : 99 sommets, 14 voisins chacun, un unique voisin commun pour deux voisins (chaque arête est dans
un seul triangle), deux voisins communs pour deux non-voisins (chaque non-arête est la diagonale d'un seul carré). Si $A$ est la matrice d'adjacence
($A_{uv}=1$ si $u\sim v$), tout cela s'écrit en une ligne :
\[
A^2+A-12I=2J\qquad(J=\text{matrice de uns}).
\]
Cette identité est \emph{le} certificat : trouver $A$ prouve l'existence ; montrer qu'aucune $A$ ne peut la satisfaire prouve la non-existence.
On sait que les valeurs propres de $A$ seraient $14,3,-4$ avec multiplicités $1,54,44$, que le groupe d'automorphismes (les symétries du graphe) est
minuscule --- $|\Aut|\in\{1,2,3\}$ (voir §2) --- et que la question est ouverte depuis 1969.

\section{Où en est la question des symétries (sources primaires)}
Une \emph{symétrie} (automorphisme) est une permutation des sommets qui envoie arêtes sur arêtes. Les résultats connus, tous vérifiés sur les textes originaux :
\begin{itemize}[nosep]
\item \textbf{Makhnev--Minakova 2004} (Diskr. Mat. 16:1, th.~1) : si $g$ est une symétrie d'ordre premier, l'ensemble $\Fix(g)$ de ses points fixes est
  un seul sommet ($|g|\in\{2,7\}$), vide ($|g|\in\{3,11\}$) ou un triangle ($|g|=3$) ; et $|\Aut|$ divise $2\cdot3^3\cdot7\cdot11$.
\item \textbf{Behbahani 2009 / Behbahani--Lam 2011} (thèse, table 21) : par énumération exhaustive de « matrices d'orbites », aucune symétrie
  d'ordre $5$, $7$ ou $11$ n'est possible, ni d'ordre $3$ avec points fixes. Il ne reste que les ordres $2$ et $3$. \textbf{L'ordre 2 n'y est pas calculé} :
  leur table commence à $p=3$ (à confirmer sur l'article de 2011).
\item \textbf{Crnković--Maksimović 2020} : pas de groupe d'ordre $6$ ni $9$. \textbf{Cesarz--Woldar 2025} : preuves sans ordinateur de résultats plus faibles.
\end{itemize}
Bilan : $|\Aut|\in\{1,2,3\}$. Un article récent (Thakkar--Severini, ar\-Xiv:2608.11211) et le dossier de départ présentaient le cas $\Z_7$ comme ouvert :
c'est une erreur de lecture, il est clos depuis 2009. Le cas d'ordre 3 est attaqué en ce moment par plusieurs projets publics.
\textbf{Le cas d'ordre 2 --- une involution --- n'a, à notre connaissance, été calculé par personne} (Keramatipour, arXiv:2604.23037, a essayé des solveurs SAT sur le problème
sans symétrie imposée, sans succès). C'est la branche que j'ai choisie : elle est, semble-t-il, vierge, et elle est rigide.

\section{Pourquoi la branche involution est rigide}
Soit $t$ une involution ($t^2=\mathrm{id}$) du graphe. Par Makhnev--Minakova, $t$ fixe exactement un sommet $x$.
\emph{Enracinons} le graphe en $x$ : ses 14 voisins forment 7 paires appariées (blocs) $B_i=\{(i,0),(i,1)\}$, car $\lambda=1$ ; chacun des 84 sommets
restants a exactement deux voisins parmi ces 14 ($\mu=2$), dans deux blocs différents ($\lambda=1$), et ces deux voisins l'identifient : les 84 sommets extérieurs
sont en bijection avec les 84 \emph{étiquettes} $\{(i,s),(j,e)\}$, $i<j$ (réduction enracinée classique, écrite notamment par Cesarz--Woldar, Keramatipour et Thakkar--Severini). Toutes les arêtes entre les 15 sommets intérieurs
et les 84 extérieurs sont donc forcées ; l'inconnue est le graphe $H$ des arêtes entre sommets extérieurs (chaque extérieur y a $12$ voisins).

Le lemme suivant est essentiellement celui de Makhnev (notes de Bielefeld : $|\Fix(t)|=1+|\Fix(t)\cap[a]|+2l$, d'où $l=0$) ; ce qui est propre à ce travail commence avec le couplage $D$.
\begin{lemma}[l'involution est forcée]
$t$ agit sur les sommets intérieurs comme le « retournement » $\sigma:(i,s)\mapsto(i,1-s)$ de chaque bloc, et sur les étiquettes par $\sigma\{a,b\}=\{a',b'\}$
(où $a'$ est le partenaire de $a$).
\end{lemma}
\begin{proof}
$t$ fixe $x$, donc préserve les 14 voisins de $x$ et leur appariement, et ne fixe aucun d'eux. Si $t$ échangeait deux sommets intérieurs $a,b$ de blocs différents,
il fixerait l'étiquette $\{a,b\}$, c'est-à-dire un sommet extérieur : contradiction. Donc chaque bloc est retourné sur lui-même.
\end{proof}

\begin{lemma}[aucun extérieur n'est voisin de son image]\label{lem:noself}
Si $u\sim tu$ alors $u$ est intérieur. En particulier $a_1(t)=14$ et les traces de $t$ sur les espaces propres $3$ et $-4$ sont nulles.
\end{lemma}
\begin{proof}
L'arête $\{u,tu\}$ est envoyée sur elle-même par $t$ ; son unique voisin commun ($\lambda=1$) est donc fixé par $t$ : c'est $x$. Donc $u\sim x$.
Avec $a_0=1$ point fixe, $a_1=14$ sommets envoyés sur un voisin, $a_2=84$ sur un non-voisin, et l'idempotent $E_3=\tfrac47I+\tfrac17A-\tfrac2{77}J$ :
$\chi_3(t)=\tfrac6{11}+14\cdot\tfrac9{77}-84\cdot\tfrac2{77}=0$, puis $\chi_{-4}(t)=a_0-1-\chi_3(t)=0$.
\end{proof}

\begin{proposition}[le couplage $D$]
Pour tout extérieur $u$, les deux voisins communs de $u$ et $\sigma u$ forment une orbite $\{v,\sigma v\}=:D(u)$, et $D$ est un \emph{couplage parfait}
sur les 42 orbites d'extérieurs. Autrement dit chaque extérieur est voisin des deux membres d'exactement une orbite.
\end{proposition}
\begin{proof}
$u\not\sim\sigma u$ (lemme~\ref{lem:noself}) donc $|N(u)\cap N(\sigma u)|=\mu=2$ ; cet ensemble est stable par $\sigma$ et sans point fixe, c'est une orbite.
Si $D(u)=\{v,\sigma v\}$ alors $u,\sigma u\in N(v)\cap N(\sigma v)$, donc $D(v)=\{u,\sigma u\}$.
\end{proof}

\paragraph{Le modèle exact.} On regroupe les 84 extérieurs en 42 orbites $U=\{u,\sigma u\}$. Pour deux orbites $U\neq V$ on note
$(a_{UV},b_{UV})=(A[u,v],A[u,\sigma v])$ : deux bits, symétriques en $U,V$. L'identité $A^2+A-12I=2J$ équivaut alors exactement à :
(R) $\sum_V(a_{UV}+b_{UV})=12$ ; (D) une seule orbite $V$ avec $a_{UV}=b_{UV}=1$ ; (I) pour chaque intérieur $c$, le nombre de voisins de $u$ dont
l'étiquette contient $c$ vaut $1$ si $c\in u$ ou $c'\in u$, et $2$ sinon ; (P) pour $U\neq V$ :
$\sum_W(a_{UW}a_{VW}+b_{UW}b_{VW})+a_{UV}=2-|u\cap v|$ et $\sum_W(a_{UW}b_{VW}+b_{UW}a_{VW})+b_{UV}=2-|u\cap\sigma v|$.
Ce sont 1722 bits et quelques milliers d'équations. C'est \emph{ce que je fais calculer} : chaque solution est un graphe de Conway (on le reconstruit et on
vérifie $A^2+A-12I=2J$ en entiers) ; l'absence de solution prouve qu'aucun graphe de Conway n'a d'involution.

\paragraph{Deux quotients.} L'involution coupe l'espace $\mathbb{R}^{99}$ en deux : les vecteurs \emph{pairs} ($f(tv)=f(v)$, dimension 50) et \emph{impairs}
($f(tv)=-f(v)$, dimension 49). Sur les pairs, $A$ agit par la \emph{matrice d'orbites} $Q$ (entrées $0,1,2$), avec $Q^2+Q-12I=2N$ ; sur les impairs par la
\emph{matrice signée} $S$ (entrées $0,\pm1$), avec $S^2+S=12I$ exactement, spectre $3^{27},(-4)^{22}$. Le bloc extérieur $S_{\mathrm{out}}$ ($42\times42$, dix
$\pm1$ par ligne) satisfait $C^{\mathsf T}S_{\mathrm{out}}=0$ et $S_{\mathrm{out}}^2+S_{\mathrm{out}}=12I-CC^{\mathsf T}$, d'où le spectre $0^7,3^{20},(-4)^{15}$
et le fait que $(S_{\mathrm{out}}+4I-\tfrac13CC^{\mathsf T})/7$ est un projecteur de rang 20. Ces structures sont ce que j'ai utilisé pour chercher des obstructions.

\section{Ce que j'ai établi : une identité de rang exacte}
L'idée vient d'un article de juin 2026 (Ishida, « No involutions in the missing Moore graph ») : regarder non pas les valeurs propres, mais les
\emph{réseaux} entiers $2$-adiques que les projecteurs découpent, et leur structure comme modules sur le groupe engendré par $t$.

\begin{theorem}
Soit $\bar\Gamma\in\F_2^{42\times42}$ la matrice d'adjacence (mod 2) du graphe des « adjacences simples » entre orbites ($a_{UV}+b_{UV}=1$) et
$N_b\in\F_2^{42\times7}$ l'incidence orbites/blocs. Alors
\[
\bar\Gamma^2+\bar\Gamma=N_bN_b^{\mathsf T}\qquad\text{et}\qquad \rg_{\F_2}\bigl(\bar\Gamma+N_bN_b^{\mathsf T}\bigr)=\rg_{\F_2}(\bar\Gamma^2)=20 .
\]
\end{theorem}
\begin{proof}
La matrice signée $S$ vérifie $S^2+S=12I$ : ses valeurs propres sont $3$ et $-4$, de multiplicités $f,g$ avec $f+g=49$ et
$3f-4g=\operatorname{tr}(S)=-7$ (bloc intérieur $-I_7$, diagonale extérieure nulle par le lemme~\ref{lem:noself}) ; donc $f=27$, $g=22$, et le polynôme
caractéristique de $S$ est $(x-3)^{27}(x+4)^{22}$.
Réduit modulo 2, il devient $(x+1)^{27}x^{22}$ ; et comme $S^2+S=12I\equiv0$, la matrice $\bar S=S\bmod2$ vérifie $\bar S^2=\bar S$, donc est diagonalisable
sur $\F_2$ : son rang est le nombre de valeurs propres égales à $1$, soit $27$. Comme $\bar S=\begin{pmatrix}I_7&N_b^{\mathsf T}\\N_b&\bar\Gamma\end{pmatrix}$,
un complément de Schur donne $\rg(\bar\Gamma+N_bN_b^{\mathsf T})=20$. L'identité $\bar\Gamma^2+\bar\Gamma=N_bN_b^{\mathsf T}$ est la réduction mod 2 de
$Q_{\mathrm{out}}^2+Q_{\mathrm{out}}=12I+4J-\mathrm{Sh}$.
\end{proof}
Ma première preuve passait par les réseaux $2$-adiques et un théorème de Broué (réseaux de $p$-permutation) ; une vérification indépendante a fait remarquer
que le spectre de $S$ et un complément de Schur suffisent. L'énoncé ne change pas, la preuve est plus courte et n'utilise que le lemme~\ref{lem:noself}.
Statut (révisé le 22 septembre) : les identités de comptage réduites mod 2 ne donnent que « $K=\bar\Gamma^2$ est idempotente de rang pair », mais la valeur $20$
découle de l'identité entière $S^2+S=12I$ et de $\operatorname{tr}S=-7$, qui sont elles-mêmes des conséquences du modèle exact. Ce n'est donc \emph{pas} une obstruction
indépendante, et ce n'est pas une contradiction (elle est compatible avec tout ce que l'on sait, y compris le plancher $\rg_{\F_7}(M)\ge27$ du dépôt public) :
c'est un invariant exact, utilisable comme filtre tardif dans une recherche. Son analogue à $k=22$ est vérifié sur le graphe $\srg(243,22,1,2)$ de Berlekamp--van Lint--Seidel.

\section{Ce que j'ai vérifié en chemin : l'audit du recensement d'ordre sept}
Le dossier recommandait d'exploiter le recensement des 208 sous-graphes induits à 7 sommets (Reimbayev, arXiv:2608.19410). J'ai transcrit les 208 formules
depuis le source \TeX{} et testé l'identité $\sum_i z_i=\binom{n}{7}$ : \textbf{elle échoue} (déficit $94\,034\,160$ à $(99,14)$). En comparant, sans les figures,
aux 208 classes déjà calculées indépendamment par le dépôt public (vague 23), exactement quatre formules sont fautives, et les corrections sont des coquilles :
$z_{21}$ (facteur $(k-3)$ en trop), $z_{80},z_{84}$ ($(2k-7)\to(2k-17)$), $z_{176}$ ($\tfrac1{84}\to\tfrac18$). Après correction, l'identité devient une identité
polynomiale en $k$ et les 208 formules coïncident avec le dépôt sur tous les coefficients. Conclusion : ce recensement n'apporte \emph{aucune} contrainte nouvelle
(seulement $2n_3\le z_{11}\le4n_3$), et la recommandation était déjà consommée par le dépôt. Détails : \texttt{notes/AUDIT\_reimbayev\_order7.md}.

\section{Ce que je fais calculer, et ce que ça a donné}
\begin{itemize}[nosep]
\item Modèles exacts (\texttt{orbit\_sat.py}, \texttt{cpsat\_orbit.py}) validés sur les contrôles : $k=4$ donne une solution qui, relevée en $9\times9$, est bien
la tour $3\times3$ ; $k=6$ est immédiatement impossible (aucun $\srg(19,6,1,2)$). Ces deux contrôles sont faibles (3 variables à $k=4$) ; contrôle fort ajouté le 22 septembre :
le graphe $\srg(243,22,1,2)$ de Berlekamp--van Lint--Seidel, avec l'involution $x\mapsto-x$, satisfait le modèle à $k=22$ sans aucune violation. Pour $k=14$ : 666\,183 variables, 1,51~M de clauses.
\item Découpage en trois cas selon l'orbite double $D(U_0)$ de l'orbite de référence (sœur, un bloc commun, disjointe) ; les autres choix s'y ramènent par symétrie.
\item Résultats : relaxation signée (CP-SAT, 25 min) : indécis. Modèle complet avec la ligne de $U_0$ entièrement fixée : indécis après $\ge30$ min de CPU
(CaDiCaL, 4 instances) et 15 min (CP-SAT). Batterie CaDiCaL/Kissat natifs (7 instances, 8~h chacune) : toutes indécises. Recuits simulés : plateaux loin de zéro.
Un dépassement de temps n'est \emph{pas} une preuve.
\end{itemize}

\section{Ce que j'en pense, et la suite}
La branche involution est entièrement formalisée, rigide, et dotée d'une structure algébrique forte (deux quotients à spectres prescrits, un couplage parfait,
un invariant de rang). Les solveurs génériques ne la décident pas en quelques heures. La méthode qui la déciderait vraiment est un \emph{générateur ordonné}
sur les 42 orbites : entrelacement spectral des sous-matrices partielles avec $\{14,3^{27},(-4)^{22}\}$, couplage $D$ et contraintes de blocs, invariant $\F_2$ en filtre,
et le groupe $2^7\rtimes S_7$ (ordre 645\,120) pour éliminer les isomorphes. C'est plusieurs jours de travail, mais l'issue serait certifiable dans les deux sens :
« aucune involution, donc $|\Aut|\in\{1,3\}$ », ou bien un graphe de Conway.

Sur la question de fond, mon intuition penche vers la non-existence (rigidité extrême de $\lambda=1,\mu=2$, absence de symétrie, plateaux des heuristiques),
mais ce n'est qu'une intuition : toutes les conditions nécessaires connues sont satisfaites, et rien dans ce travail ne tranche. Les calculs sont codés
symétriquement et rendront ce qu'ils rendront.
\end{document}
